My plan to advance the frontier of gravitational wave cosmology, using both instrumental and observational techniques, is as follows. 
First, 
I will explore new astronomical instruments designed to complement \Rubin\ using spectroscopy, 
while also deploying monitoring tools for the \LSSTCam\ to improve \LSST\ data quality. 
Second, I will lead \ToO\ searches to detect the next bright siren and other exotic astrophysical events. 
Third, I will usher standard siren cosmology into a precision era, leveraging \LSST\ data alongside specialized observational programs for bright and dark siren cosmology. 
Together, these efforts span instrumentation, transient discovery, and precision cosmological measurement, all with the goal of gaining a better understanding of dark energy.

I am committed to crossing disciplinary boundaries to ensure that scientific results reflect perspectives from all forms of expertise. The field of gravitational wave cosmology is inherently interdisciplinary, leveraging specialized knowledge from both gravitational physics and cosmological science, in addition to incorporating recent performance improvements in parameter inference driven by the adoption of GPUs for analysis pipelines. \Institute\ provides an exquisite base of knowledge, and I will work with members of the department of physics, across theoretical, experimental, and observational expertise to develop standard siren cosmology into an era of precision measurement. In addition to my host department, I will collaborate with the computer science department to revise our analysis techniques to leverage the latest hardware and computational methods. This will be especially important as the quality of galaxy catalog data is set to grow with the onset of the \LSST\ in mid-2026, and as samples of \GW s will increase with improved detector sensitivities as early as 2029.

I am eager to act as a mentor for other early-career scientists who are eager to explore the cosmic frontier. This passion for mentorship has grown throughout my life and career, as one of the most rewarding experiences I felt in science was when I led high school students from my hometown in guided lectures and observations of the 2024 Solar Eclipse from the path of totality. I will continue engaging with the next generation of burgeoning scientists at both the PhD and undergraduate level, mentoring  physics PhD students as they take on their rotations every quarter, hoping to guide them toward the research field of their choice, whether it is in cosmology or some other field. In addition to supporting PhD level students, I would like to serve as a mentor for SULI summer students at KIPAC, providing undergraduate students a supportive mentor as many encounter university-level research for the first time.

I am keen to engage with a unique interdisciplinary environment where emerging ideas are developed into innovative methods using different forms of expertise. I believe my previous research experience, future research plans, and desire to collaborate across disciplines and experience levels will contribute to a productive and innovative research program as a \Position. 